% --- Archon physics formalization source begin ---
% archon:physics
% archon:covers PhyXMiniProblems/problem_phyx_mini_0146.lean
% archon:source-report reports/phyx_mini/problem_phyx_mini_0146.source.json

\chapter{Physics problem phyx\_mini\_0146}
\label{ch:PhyXMiniProblems_problem_phyx_mini_0146}

\paragraph{Problem source.}
\#\# Physical scenario

Light is incident on an equilateral glass prism at a \textbackslash{}(45.0\textasciicircum{}\textbackslash{}circ\textbackslash{}) angle to one face, figure. Assume that \textbackslash{}(n = 1.54\textbackslash{}).

\#\# Question

Calculate the angle at which light emerges from the opposite face.

\#\# Answer choices

- A: A: \$66.2\textasciicircum{}\{\textbackslash{}circ\}\textbackslash{}

- B: B: \$46.2\textasciicircum{}\{\textbackslash{}circ\}\textbackslash{}

- C: C: \$56.2\textasciicircum{}\{\textbackslash{}circ\}\textbackslash{}

- D: D: \$76.2\textasciicircum{}\{\textbackslash{}circ\}\textbackslash{}

\#\# Figure caption (auxiliary; use the image as primary evidence)

The image depicts a diagram involving a prism, light rays, and angles. Here's a detailed description:

- **Prism**: At the center, there's a triangular prism shaded in light green.
- **Incident Ray**: A pink arrow enters the prism from the left side at an angle of 45.0° with respect to the normal at the entry point.
- **Refracted Ray**: Inside the prism, the ray bends towards the right side of the prism.
- **Emergent Ray**: The refracted ray exits the prism and bends again, forming an angle with the normal at the exit point. This angle is labeled with a question mark (?).
- **Dashed Lines**: There are dashed lines representing normals at the points where the light ray enters and exits the prism, extending outside the prism surface.
- **Angles**: The angle on the left is labeled as 45.0°, while the angle on the right is marked with a question mark.

The diagram illustrates the refraction of light through a prism, focusing on determining the unknown angle of refraction at the exit point.

\#\# Dataset metadata

Geometrical Optics; Physical Model Grounding Reasoning, Spatial Relation Reasoning

\paragraph{Recorded answer/context.}
C: C: \$56.2\textasciicircum{}\{\textbackslash{}circ\}\textbackslash{}

\paragraph{Figure/image path.}
/root/proposal\_for\_physic/science-mango/phyx\_mini\_run/phyx\_data/test\_image/146.png

\paragraph{Formalization target.}
create a compiling Lean file with sorry bodies at `PhyXMiniProblems/problem\_phyx\_mini\_0146.lean`. The Lean declarations must preserve the physical quantities, dimensions or dimensional roles, figure labels, governing-law hypotheses, and final relation expressed by this problem.
Use Mathlib/Physlib names found through LeanExplore where available. If a physics API is missing, introduce faithful local abstractions rather than scalar placeholder aliases.

\begin{theorem}[Physics formalization target]
\label{thm:physics:phyx_mini_0146:target}
\lean{PhyXMiniProblems.ProblemPhyXMini0146.emergenceAngleIsChoiceC}
\uses{def:physics:phyx-mini-0146:phyxminiproblems-problemphyxmini0146-equilateralglassprismdiagram, def:physics:phyx-mini-0146:phyxminiproblems-problemphyxmini0146-matchesproblemdata, def:physics:phyx-mini-0146:phyxminiproblems-problemphyxmini0146-hasphysicalopticalparameters, def:physics:phyx-mini-0146:phyxminiproblems-problemphyxmini0146-satisfiesprismraygeometry, def:physics:phyx-mini-0146:phyxminiproblems-problemphyxmini0146-obeyssnellslaw, def:physics:phyx-mini-0146:phyxminiproblems-problemphyxmini0146-answerchoice, def:physics:phyx-mini-0146:phyxminiproblems-problemphyxmini0146-roundstonearesttenthdegree, def:physics:phyx-mini-0146:phyxminiproblems-problemphyxmini0146-isuniquenearestanswerchoice}
The assigned autoformalize agent should translate this physics problem into Lean declarations in the covered file, with theorem and lemma proof bodies written as `by sorry`.
\end{theorem}
\begin{proof}
This is an autoformalization task, not a proof task. Produce statements that can later be proved by the physics prover without weakening the physical model.
\end{proof}

% --- Archon declaration topology begin ---
\section{Declaration topology}
\begin{definition}[angle From Degrees]
  \label{def:physics:phyx-mini-0146:phyxminiproblems-problemphyxmini0146-anglefromdegrees}
  \lean{PhyXMiniProblems.ProblemPhyXMini0146.angleFromDegrees}
  Convert a scalar degree readout to a physical angle
\end{definition}
\begin{proof}
  This is the declared model definition of angle From Degrees.
\end{proof}
\begin{definition}[angle Degrees Readout]
  \label{def:physics:phyx-mini-0146:phyxminiproblems-problemphyxmini0146-angledegreesreadout}
  \lean{PhyXMiniProblems.ProblemPhyXMini0146.angleDegreesReadout}
  Principal degree readout of a physical angle
\end{definition}
\begin{proof}
  This is the declared model definition of angle Degrees Readout.
\end{proof}
\begin{definition}[Optical Medium]
  \label{def:physics:phyx-mini-0146:phyxminiproblems-problemphyxmini0146-opticalmedium}
  \lean{PhyXMiniProblems.ProblemPhyXMini0146.OpticalMedium}
  The two homogeneous optical media traversed by the ray
\end{definition}
\begin{proof}
  This is the declared model definition of Optical Medium.
\end{proof}
\begin{definition}[Prism Interface]
  \label{def:physics:phyx-mini-0146:phyxminiproblems-problemphyxmini0146-prisminterface}
  \lean{PhyXMiniProblems.ProblemPhyXMini0146.PrismInterface}
  The two sloping prism faces crossed by the light ray
\end{definition}
\begin{proof}
  This is the declared model definition of Prism Interface.
\end{proof}
\begin{definition}[Ray Segment]
  \label{def:physics:phyx-mini-0146:phyxminiproblems-problemphyxmini0146-raysegment}
  \lean{PhyXMiniProblems.ProblemPhyXMini0146.RaySegment}
  The three directed portions of the light path visible in the image
\end{definition}
\begin{proof}
  This is the declared model definition of Ray Segment.
\end{proof}
\begin{definition}[Figure Angle Label]
  \label{def:physics:phyx-mini-0146:phyxminiproblems-problemphyxmini0146-figureanglelabel}
  \lean{PhyXMiniProblems.ProblemPhyXMini0146.FigureAngleLabel}
  The two normal-referenced angle labels in the primary image
\end{definition}
\begin{proof}
  This is the declared model definition of Figure Angle Label.
\end{proof}
\begin{definition}[incident Medium]
  \label{def:physics:phyx-mini-0146:phyxminiproblems-problemphyxmini0146-incidentmedium}
  \lean{PhyXMiniProblems.ProblemPhyXMini0146.incidentMedium}
  \uses{def:physics:phyx-mini-0146:phyxminiproblems-problemphyxmini0146-opticalmedium, def:physics:phyx-mini-0146:phyxminiproblems-problemphyxmini0146-prisminterface}
  Medium containing the ray immediately before each interface crossing
\end{definition}
\begin{proof}
  This is the declared model definition of incident Medium.
\end{proof}
\begin{definition}[transmitted Medium]
  \label{def:physics:phyx-mini-0146:phyxminiproblems-problemphyxmini0146-transmittedmedium}
  \lean{PhyXMiniProblems.ProblemPhyXMini0146.transmittedMedium}
  \uses{def:physics:phyx-mini-0146:phyxminiproblems-problemphyxmini0146-opticalmedium, def:physics:phyx-mini-0146:phyxminiproblems-problemphyxmini0146-prisminterface}
  Medium containing the ray immediately after each interface crossing
\end{definition}
\begin{proof}
  This is the declared model definition of transmitted Medium.
\end{proof}
\begin{definition}[Equilateral Glass Prism Diagram]
  \label{def:physics:phyx-mini-0146:phyxminiproblems-problemphyxmini0146-equilateralglassprismdiagram}
  \lean{PhyXMiniProblems.ProblemPhyXMini0146.EquilateralGlassPrismDiagram}
  \uses{def:physics:phyx-mini-0146:phyxminiproblems-problemphyxmini0146-opticalmedium, def:physics:phyx-mini-0146:phyxminiproblems-problemphyxmini0146-prisminterface, def:physics:phyx-mini-0146:phyxminiproblems-problemphyxmini0146-raysegment, def:physics:phyx-mini-0146:phyxminiproblems-problemphyxmini0146-figureanglelabel}
  The physical quantities attached to the labelled equilateral-prism diagram. At each face, incidenceAngle is measured from the dashed normal on the incident side and refractionAngle from that normal on the transmitted side. Thus the exit-face refraction angle is the emergence angle requested by the problem
\end{definition}
\begin{proof}
  This is the declared model definition of Equilateral Glass Prism Diagram.
\end{proof}
\begin{definition}[Matches Problem Data]
  \label{def:physics:phyx-mini-0146:phyxminiproblems-problemphyxmini0146-matchesproblemdata}
  \lean{PhyXMiniProblems.ProblemPhyXMini0146.MatchesProblemData}
  \uses{def:physics:phyx-mini-0146:phyxminiproblems-problemphyxmini0146-anglefromdegrees, def:physics:phyx-mini-0146:phyxminiproblems-problemphyxmini0146-equilateralglassprismdiagram}
  Primary-image readouts and numerical data stated in the problem. In particular, this does not assign a numerical value to the exit question mark
\end{definition}
\begin{proof}
  This is the declared model definition of Matches Problem Data.
\end{proof}
\begin{definition}[Is Physical Ray Angle]
  \label{def:physics:phyx-mini-0146:phyxminiproblems-problemphyxmini0146-isphysicalrayangle}
  \lean{PhyXMiniProblems.ProblemPhyXMini0146.IsPhysicalRayAngle}
  An angle on the acute principal branch used by the depicted ray
\end{definition}
\begin{proof}
  This is the declared model definition of Is Physical Ray Angle.
\end{proof}
\begin{definition}[Has Physical Optical Parameters]
  \label{def:physics:phyx-mini-0146:phyxminiproblems-problemphyxmini0146-hasphysicalopticalparameters}
  \lean{PhyXMiniProblems.ProblemPhyXMini0146.HasPhysicalOpticalParameters}
  \uses{def:physics:phyx-mini-0146:phyxminiproblems-problemphyxmini0146-equilateralglassprismdiagram, def:physics:phyx-mini-0146:phyxminiproblems-problemphyxmini0146-isphysicalrayangle}
  Positivity and principal-branch conditions for the optical quantities
\end{definition}
\begin{proof}
  This is the declared model definition of Has Physical Optical Parameters.
\end{proof}
\begin{definition}[Satisfies Prism Ray Geometry]
  \label{def:physics:phyx-mini-0146:phyxminiproblems-problemphyxmini0146-satisfiesprismraygeometry}
  \lean{PhyXMiniProblems.ProblemPhyXMini0146.SatisfiesPrismRayGeometry}
  \uses{def:physics:phyx-mini-0146:phyxminiproblems-problemphyxmini0146-equilateralglassprismdiagram}
  Standard prism geometry for the internal ray: its refraction angle at entry and incidence angle at exit add to the angle between the refracting faces
\end{definition}
\begin{proof}
  This is the declared model definition of Satisfies Prism Ray Geometry.
\end{proof}
\begin{definition}[Satisfies Snells Law At]
  \label{def:physics:phyx-mini-0146:phyxminiproblems-problemphyxmini0146-satisfiessnellslawat}
  \lean{PhyXMiniProblems.ProblemPhyXMini0146.SatisfiesSnellsLawAt}
  \uses{def:physics:phyx-mini-0146:phyxminiproblems-problemphyxmini0146-prisminterface, def:physics:phyx-mini-0146:phyxminiproblems-problemphyxmini0146-incidentmedium, def:physics:phyx-mini-0146:phyxminiproblems-problemphyxmini0146-transmittedmedium, def:physics:phyx-mini-0146:phyxminiproblems-problemphyxmini0146-equilateralglassprismdiagram}
  Snell's law n₁ sin θ₁ = n₂ sin θ₂ at one crossed prism face
\end{definition}
\begin{proof}
  This is the declared model definition of Satisfies Snells Law At.
\end{proof}
\begin{definition}[Obeys Snells Law]
  \label{def:physics:phyx-mini-0146:phyxminiproblems-problemphyxmini0146-obeyssnellslaw}
  \lean{PhyXMiniProblems.ProblemPhyXMini0146.ObeysSnellsLaw}
  \uses{def:physics:phyx-mini-0146:phyxminiproblems-problemphyxmini0146-equilateralglassprismdiagram, def:physics:phyx-mini-0146:phyxminiproblems-problemphyxmini0146-satisfiessnellslawat}
  The depicted ray obeys Snell's law at both prism faces
\end{definition}
\begin{proof}
  This is the declared model definition of Obeys Snells Law.
\end{proof}
\begin{definition}[Answer Choice]
  \label{def:physics:phyx-mini-0146:phyxminiproblems-problemphyxmini0146-answerchoice}
  \lean{PhyXMiniProblems.ProblemPhyXMini0146.AnswerChoice}
  Labels of the four displayed answer choices
\end{definition}
\begin{proof}
  This is the declared model definition of Answer Choice.
\end{proof}
\begin{definition}[angle Degrees]
  \label{def:physics:phyx-mini-0146:phyxminiproblems-problemphyxmini0146-answerchoice-angledegrees}
  \lean{PhyXMiniProblems.ProblemPhyXMini0146.AnswerChoice.angleDegrees}
  \uses{def:physics:phyx-mini-0146:phyxminiproblems-problemphyxmini0146-answerchoice}
  Emergence-angle readout in degrees printed beside each answer choice
\end{definition}
\begin{proof}
  This is the declared model definition of angle Degrees.
\end{proof}
\begin{definition}[Rounds To Nearest Tenth Degree]
  \label{def:physics:phyx-mini-0146:phyxminiproblems-problemphyxmini0146-roundstonearesttenthdegree}
  \lean{PhyXMiniProblems.ProblemPhyXMini0146.RoundsToNearestTenthDegree}
  \uses{def:physics:phyx-mini-0146:phyxminiproblems-problemphyxmini0146-angledegreesreadout}
  A physical angle rounds to a displayed degree value to the nearest tenth
\end{definition}
\begin{proof}
  This is the declared model definition of Rounds To Nearest Tenth Degree.
\end{proof}
\begin{definition}[Is Unique Nearest Answer Choice]
  \label{def:physics:phyx-mini-0146:phyxminiproblems-problemphyxmini0146-isuniquenearestanswerchoice}
  \lean{PhyXMiniProblems.ProblemPhyXMini0146.IsUniqueNearestAnswerChoice}
  \uses{def:physics:phyx-mini-0146:phyxminiproblems-problemphyxmini0146-angledegreesreadout, def:physics:phyx-mini-0146:phyxminiproblems-problemphyxmini0146-answerchoice}
  A choice is strictly nearer to the physical angle than every other choice
\end{definition}
\begin{proof}
  This is the declared model definition of Is Unique Nearest Answer Choice.
\end{proof}
% --- Archon declaration topology end ---
% --- Archon physics formalization source end ---
